质量矩阵与夸克带电流的基失配都已确定后，费米子相互作用可以统一写到质量本征基。对最小可重整化标准模型，电磁流与中性流在味空间对角，而夸克带电流含 $V_{\rm CKM}$：
\begin{align}
 \Lag_{\rm int}^{f}
 ={}&-\hbar c\,g_{\rm em}\sum_fQ_f\bar f\gamma^\mu f\,A_\mu^{\rm geo}
 -\frac{\hbar c\,g}{2c_W}\sum_f
 \bar f\gamma^\mu(g_V^f-g_A^f\gamma^5)f\,\mathcal Z_\mu
 \nonumber\\
 &-\frac{\hbar c\,g}{\sqrt2}
 \left[\bar u_i\gamma^\mu P_L(V_{\rm CKM})_{ij}d_j\,\mathcal W_\mu^+
 +\bar d_j\gamma^\mu P_L(V_{\rm CKM})_{ij}^*u_i\,\mathcal W_\mu^-\right].
 \label{eq:sm-mass-basis-fermion-interactions-dp11}
\end{align}
其中 $f$ 遍历带电费米子与中微子中性流所需的质量本征场；严格最小标准模型中的无质量中微子基可与带电轻子基对齐，因此不产生可观测轻子混合矩阵。对全部外腿采用流入动量约定，相应的约化费米子顶角为
\begin{align}
 \gamma f\bar f:&\qquad -\ii g_{\rm em}Q_f\gamma^\mu,\\
 Zf\bar f:&\qquad -\ii\frac{g}{2c_W}\gamma^\mu(g_V^f-g_A^f\gamma^5),\\
 W^+u_i\bar d_j:&\qquad -\ii\frac{g}{\sqrt2}(V_{\rm CKM})_{ij}\gamma^\mu P_L.
 \label{eq:sm-fermion-feynman-rules-dp11}
\end{align}
电磁与中性流的味对角性来自世代空间中的单位算符；带电流中的 $V_{\rm CKM}$ 则来自两套左手 Yukawa 奇异矢的失配。
